\documentclass{article}
\begin{document}
\title{The placement of old effects hidden at the last cut matters}
\maketitle

\paragraph{An irregular receipt law.}
Keep the same exclusive clipped brake plant, speed
\(v\), no-loss world and at-most-one omission, with
deadline \(2v+1\). A publicly known set \(U\)
of nominal slots would reveal their omission only
after the deadline. All other useful omissions are
reported before the next action. In both classes
the plant does not reveal unlabeled hidden speed.
This is deliberately different from a common
receipt delay: the profile can be altered by which
receiver is active during particular parts of the old
plan. Here that profile is stipulated, not a
hardware source fact.

With the gap limit removed, a winning policy exists
exactly when \(|U|\le v\). The no-loss prefix is balanced,
with \(v\) assignments to each actuator. The final action
has to remedy any one unreported missing slot, so all
of \(U\) must have the same actuator assignment in that
prefix. Conversely choose any balanced word with that
property. Use it before a qualified receipt, insert
a timely repair when named, and otherwise correct the
common \(\sigma\) sign at the final cut (clips in the
no-loss world). The no-loss progress argument is
unchanged by which sign is chosen. A uniform use of the
maximum possible receipt delay would wrongly rule out
some of these progress cases.

\paragraph{Exact reduction for checking one fixed word.}
For the still-unacknowledged worlds of a fixed winning word,
a lost grant of sign \(\sigma\) changes only one robot's
speed by one unit until its eventual final compensating
grant. For unit slots the displacement difference accumulates
monotonically, with a half-ramp at the omitted slot and
a ramp removed at compensation. If two possible omissions
have the same sign, the earlier loss produces a pointwise
larger deviation in that direction. All still-unacknowledged
trajectories are therefore between the no-loss trajectory
and the earliest possible unacknowledged omission, even
during the final repair. This is a sufficiency kernel for
one fixed word and this one-loss timing law; it does not
erase the other indices when finding a legal word, or the
different branches created by earlier named receipts.

\paragraph{Same count and maximum delay, different slack order.}
For \(v\ge2\) let one opaque set be
\(\{0\}\cup\{v+1,\ldots,2v-1\}\).
The no-loss word is necessarily \(-,+^{v},-^{v-1}\)
up to reversal of signs. Erase slot zero. No receipt arrives
in time; the effective sequence is all \(v\) of one sign
then \(v\) of the other, with final displacement magnitude
\(v^2\). Conversely that word and the repair above
fit \(G=v^2\): a timely loss in the middle block adds
only one coast at speed of magnitude at most \(v-2\)
(or one at speed \(-1\)), a late tail loss adds at most
\(v-1\) to the nominal terminal \(v^2-2v\), and the
slot-zero bound is attained.

For the other opaque set take all positive slots of
the nominal alternating ordered-pair word
\(+, -,-,+,\ldots\), with a final ordered pair
if needed. This too has \(v\) opaque slots, including
zero, and the same maximum latency if every opaque
receipt comes just after the last cut. Its nominal gap
oscillates from 0 to 1. A timely repair shifts it
by at most 1, and an opaque omission of the same
positive sign changes it by at most \(2v\) in the negative
direction. Therefore this profile needs at most \(2v\)
symmetric slack, versus a sharp \(v^2\) for the first.
The two summaries (number and maximum delay) cannot
support a sharp geometric admission or reuse decision.
This difference is between evidence placements, not
performance of two policies granted different physical
actuators. Retain the actual law relating receipt times
to which previous effects they qualify.

The constructive replay checks the pointwise hull and
exact small-family feasibility across all opaque masks.
No real receipt source, independent mission, or claim
against split convex actuation is included.
\end{document}
